\documentclass[11pt]{article}

\usepackage[a4paper,margin=1in]{geometry}
\usepackage[T1]{fontenc}
\usepackage[utf8]{inputenc}
\usepackage{lmodern}
\usepackage{amsmath,amssymb,mathtools}
\usepackage{booktabs}
\usepackage{hyperref}
\usepackage{xcolor}

\hypersetup{
  colorlinks=true,
  linkcolor=blue!60!black,
  citecolor=blue!60!black,
  urlcolor=blue!60!black
}

\newcommand{\Ds}{D_s}
\newcommand{\Dsone}{D_{s1}}
\newcommand{\ii}{\mathrm{i}}
\newcommand{\ee}{\mathrm{e}}
\newcommand{\Tr}{\mathrm{Tr}}
\newcommand{\slashp}[1]{#1\!\!\!/}

\title{Pre-Automode Report: $\Dsone\to\Ds\gamma$ LCSR Setup}
\author{Ds gamma project}
\date{\today}

\begin{document}
\maketitle

\section{Purpose}

This report summarizes the choices fixed before the autonomous calculation of
the radiative transitions
\[
  \Dsone(2460)\to\Ds\gamma,\qquad
  \Dsone(2536)\to\Ds\gamma .
\]
The goal of the next phase is to compute the form factor $G$ and the decay
width
\[
  \Gamma(1^+\to 0^-\gamma)=\frac{\alpha_{\rm em}}{3}\,G^2\,|\vec q_\gamma|^3,
  \qquad
  |\vec q_\gamma|
  =
  \frac{m_{\Dsone}^2-m_{\Ds}^2}{2m_{\Dsone}} .
\]

\section{Fixed Conventions}

We use the mostly-minus metric, $g_{\mu\nu}=\mathrm{diag}(+1,-1,-1,-1)$, and
\[
  \sigma_{\mu\nu}=\frac{\ii}{2}[\gamma_\mu,\gamma_\nu],
  \qquad
  \Tr[\gamma_5\gamma_\mu\gamma_\nu\gamma_\alpha\gamma_\beta]
  =
  -4\ii\,\epsilon_{\mu\nu\alpha\beta}.
\]
The momentum assignment is
\[
  P=p+q,
\]
where $P$ is the initial $\Dsone$ momentum, $p$ is the final $\Ds$ momentum, and
$q$ is the photon momentum.

The interpolating currents are
\begin{align}
  J^{(0)}_{A\mu}(x)
  &= \bar s(x)\gamma_\mu\gamma_5 Q(x),\\
  J^{(0)}_{B\mu}(x)
  &= \frac{\ii}{m_Q+m_s}\,
     \bar s(x)\sigma_{\mu\alpha}P^\alpha\gamma_5Q(x),\\
  J_{\Ds}(x)
  &= \bar s(x)\ii\gamma_5Q(x).
\end{align}
The tensor current contracts with $P^\alpha$, not with the final-state momentum
$p^\alpha$.

\section{OPE Organization}

The OPE is organized as
\[
  \Pi_\mu^{\rm OPE}
  =
  \Pi_\mu^{\rm hard}
  +
  \Pi_\mu^{\rm soft}.
\]

\subsection{Hard photon emission}

The hard sector contains photon emission from both quark lines:
\[
  \Pi_\mu^{\rm hard}
  =
  \ii N_c\int d^4x\,\ee^{\ii p\cdot x}
  \Tr_D\!\left[
    \Gamma_\mu S_Q^\gamma(x)\ii\gamma_5S_s^0(-x)
    +
    \Gamma_\mu S_Q^0(x)\ii\gamma_5S_s^\gamma(-x)
  \right].
\]
Thus both $S_Q^\gamma$ and $S_s^\gamma$ are included.  The generic
electromagnetic background-field correction is
\[
  S_q^\gamma(x)
  =
  -\ii e_q
  \int\frac{d^4k}{(2\pi)^4}\,\ee^{-\ii k\cdot x}
  \int_0^1du
  \left[
    \frac{\slashp{k}+m_q}{2(m_q^2-k^2)^2}
    F^{\alpha\beta}(ux)\sigma_{\alpha\beta}
    +
    \frac{u x_\alpha}{m_q^2-k^2}
    F^{\alpha\beta}(ux)\gamma_\beta
  \right].
\]
It is used with $(q,m_q,e_q)=(Q,m_Q,e_Q)$ for the heavy line and
$(s,m_s,e_s)$ for the strange line.

The base calculation does not use ordinary condensate-expanded quark
propagators.  In particular, terms such as
$\langle\bar s s\rangle$, $\langle\bar s g_s\sigma Gs\rangle$, and
$\langle G^2\rangle$ are not inserted as propagator corrections.

\subsection{Soft photon emission}

The soft sector uses photon distribution amplitudes (DAs).  For two-particle
DAs,
\[
  \Pi_\mu^{\rm soft,2p}
  =
  -\frac{\ii}{4}\int d^4x\,\ee^{\ii p\cdot x}
  \sum_i
  \Tr_D[\Gamma_\mu S_Q^0(x)\ii\gamma_5\Gamma_i]\,
  \langle\gamma(q,\varepsilon)|
  \bar s(x)\Gamma_i s(0)
  |0\rangle .
\]
Condensate parameters that appear inside photon DA normalizations, such as
$\chi\langle\bar s s\rangle$, are kept as photon-wavefunction inputs.  They are
not counted as condensate-expanded propagators.

For three-particle DAs, the background-gluon propagator insertion is needed:
\[
  S_Q^G(x)
  =
  -\ii g_s
  \int\frac{d^4k}{(2\pi)^4}\,\ee^{-\ii k\cdot x}
  \int_0^1du
  \left[
    \frac{\slashp{k}+m_Q}{2(m_Q^2-k^2)^2}
    G^{\alpha\beta}(ux)\sigma_{\alpha\beta}
    +
    \frac{u x_\alpha}{m_Q^2-k^2}
    G^{\alpha\beta}(ux)\gamma_\beta
  \right].
\]
This is not a gluon-condensate correction; it is the operator insertion needed
for matrix elements of the form
$\langle\gamma|\bar s\Gamma g_sG_{\alpha\beta}s|0\rangle$.

\section{Staged Plan}

\begin{description}
  \item[Stage 1] Compute a controlled baseline including hard photon emission
  and two-particle photon DAs.  Keep all explicit strange-quark mass terms,
  including possible $m_s^2$ contributions.

  \item[Stage 2] Add the three-particle photon DA terms generated by the
  background-gluon insertion $S_Q^G$.
\end{description}

\section{Numerical Inputs Fixed}

\begin{center}
\begin{tabular}{lll}
\toprule
Input & Choice & Comment\\
\midrule
$\chi(1\,{\rm GeV})$ & $+3.15\pm0.30\,{\rm GeV}^{-2}$ & sign fixed before automode\\
$\theta$ & $35.3^\circ$ central & scan after central run\\
$m_s$ & retained & include explicit $m_s$ and $m_s^2$ terms\\
Photon DAs & \texttt{photon\_da.py} & through twist 4\\
\bottomrule
\end{tabular}
\end{center}

\section{Implemented Artifacts}

\begin{itemize}
  \item \texttt{conventions.md}: core conventions and current definitions.
  \item \texttt{inputs\_table.csv}: numerical inputs, now with positive
  $\chi=+3.15\,{\rm GeV}^{-2}$.
  \item \texttt{photon\_da.py}: photon DA functions for Stage 1 and Stage 2.
  \item \texttt{scripts/step2\_current\_building\_blocks.wl}: FeynCalc checks
  of Dirac conventions and current vertices.
  \item \texttt{scripts/step2\_hard\_photon\_numerators.wl}: numerator-level
  hard-photon traces for heavy-line and strange-line emission.
  \item \texttt{notes/Ds1\_to\_Ds\_gamma\_LCSR\_notes.tex}: living derivation
  notes compiled to PDF.
\end{itemize}

\section{Next Automode Tasks}

\begin{enumerate}
  \item Derive the Stage 1 invariant amplitude for the chosen E1 Lorentz
  structure.
  \item Implement the Borel-transformed Stage 1 sum rule.
  \item Determine reasonable Borel windows and continuum thresholds.
  \item Extract central $G_A$ and $G_B$, then form the physical mixed
  amplitudes.
  \item Compute $\Gamma[\Dsone(2460)\to\Ds\gamma]$ and
  $\Gamma[\Dsone(2536)\to\Ds\gamma]$.
  \item Scan the mixing angle $\theta$ and selected input uncertainties.
  \item Add Stage 2 as a controlled correction.
\end{enumerate}

\end{document}
